\chapter{Kontrast}
Kontrast to ważny temat dla naszej aplikacji, ponieważ cichlibyśmy żeby była ona dostępna dla jak największej ilości osób. Dlatego staramy się żeby tekst wyświetlany na ekranie był jak najlepiej widoczny.
\section{Ogólny kontrast}
Tła wewnątrz aplikacji maję 3 główne wartości (rysunek \ref{fig:bgpalet})

\begin{itemize}
    \item \textbf{Najciemniejszy} \#1e1e1e
    \item \textbf{Średnio ciemny} \#2f2f2f
    \item \textbf{Najjaśniejszy} \#585858
\end{itemize}

\MyFigure[0.5\textwidth]{images/kontrast/palataKolorówTła.png}{Paleta kolorów tła}{fig:bgpalet}{Ręcznie dobrane kolory}

Biały tekst na najjaśniejszym z tych kolorów posiada wartość 7.11 (rysunek \ref{fig:brightest}). Co spełnia zarówna wymagania AA jak i AAA.

\MyFigure[0.5\textwidth]{images/kontrast/brightest.png}{Biały napis na najjaśniejszym kolorze tła}{fig:brightest}{\url{https://colourcontrast.cc/?background=585858\&foreground=ffffff}}
Na ekranie głównym oraz w paru innych miejscach możemy róznież za obserwować tekst oraz guziki koloru niebieskiego (\#4b81ff). Ten kolor pojawia się jedynie na najciemniejszym tle a jego kontrast w tej sytuacji wynosi 4.67 (rysunek \ref{fig:bluetest}) co jest dopuszczalne w AA oraz AAA dla dużego tekstu.
\MyFigure[0.5\textwidth]{images/kontrast/bluetest.png}{Niebieski na najciemniejszym kolorze tła}{fig:bluetest}{\url{https://colourcontrast.cc/?background=585858\&foreground=ffffff}}

Ostatnią kombinacją kolorów w naszej aplikacji są niebierskie przyciuski z białymi napisami lub symbolami. Głównie w jumper'ach oraz na ekranie do zarządzania git'em.
\MyFigure[0.5\textwidth]{images/elementy/bluebutton.png}{Niebieski guzik na bialym}{fig:bluebutton}{\url{https://colourcontrast.cc/?background=325fc7\&foreground=ffffff}}

Ich kontrast wynosi 5.84 więc również jest git dla użych napisóœ zarówno w AAA jak i AA.
\MyFigure[0.5\textwidth]{images/kontrast/bluebutton.png}{Biały napis na niebieskim tle}{fig:bluebuttonk}{Ręcznie dobrane kolory}
\section{Kontrast Kodu}
Nasza aplikacja oferuje 4 schematy kolorów w edytorze kodu wszystkie prezentowane na tle \#1e1e1e. Schematy można wybierać z poziomu ustawień aplikacji. Każdy schemat ma swoje zastosowania i może być użyty przez różnych ludzi.
\subsection{Domyślny}
\MyFigure[0.5\textwidth]{images/schematy/default.png}{Schemat koloru domyślny}{fig:default}{\parencite{makieta}}
Podstawowy schemat kolorów inspirowany klasycznym motywem Dark+ w odróżnieniu do większości współczesnych motywów jest z tonowany i elegancki. Przypomina mi czasy kiedy zaczynałem moją przygodę z programowaniem.

\begin{table}[ht]
\centering
\caption{Schemat kolorów domyślny}
\label{tab:default}
\begin{tabular}{|>{\raggedright\arraybackslash}p{4cm}|c|c|}
\hline
\textbf{Element} & \textbf{Kolor Hex} & \textbf{Kontrast} \\
\hline
funkcja       & \#D8DD73 & 11.52:1 \\
operator      & \#BFBFAF & 8.96:1 \\
słowo kluczowe & \#E161D4 & 5.45:1 \\
zmienna       & \#8FC9EA & 9.3:1 \\
klasa         & \#48B988 & 6.8:1 \\
typ           & \#4795D9 & 5.21:1 \\
string        & \#C88B72 & 5.87:1 \\
stała         & \#9CDBF0 & 10.96:1 \\
liczba        & \#AAC295 & 8.62:1 \\
numer linii   & \#59B4DE & 7.15:1 \\
\hline
ostrzeżenie   & \#ECA25C & 7.85:1 \\
error         & \#ED9292 & 7.24:1 \\
\hline
\end{tabular}
\newline\newline
{\raggedright\fontsize{10pt}{12pt}\selectfont\textit{Źródło: Opracowanie własne}\par}
\end{table}

\subsection{Drugi}
\MyFigure[0.5\textwidth]{images/schematy/second.png}{Schemat koloru drugi}{fig:secong}{\parencite{makieta}}
Drugi motyw jest nieco bardziej edżi dla buntowników, emo lub ludzi bez smaku. Jakikolwiek balans jest wyrzucony za okno w imie jaskrawych kolorów bez żadnej klasy. Ludzie używający tego motywu nie wiedzieli by co to dobry motyw gdyby usiadł im na twarzy.

\begin{table}[ht]
\centering
\caption{Schemat kolorów drugi}
\label{tab:alt1}
\begin{tabular}{|>{\raggedright\arraybackslash}p{4cm}|c|c|}
\hline
\textbf{Element} & \textbf{Kolor Hex} & \textbf{Kontrast} \\
\hline
funkcja       & \#56B4E9 & 7.22:1 \\
operator      & \#CCCCCC & 10.38:1 \\
słowo kluczowe & \#E69F00 & 7.4:1 \\
zmienna       & \#56B4E9 & 7.22:1 \\
klasa         & \#009E73 & 4.87:1 \\
typ           & \#E69F00 & 7.4:1 \\
string        & \#F0E442 & 12.61:1 \\
stała         & \#F0E442 & 12.61:1 \\
liczba        & \#CC79A7 & 5.45:1 \\
numer linii   & \#59B4DE & 7.15:1 \\
\hline
ostrzeżenie   & \#ECA25C & 7.85:1 \\
error         & \#ED9292 & 7.24:1 \\
\hline
\end{tabular}
\newline\newline
{\raggedright\fontsize{10pt}{12pt}\selectfont\textit{Źródło: Opracowanie własne}\par}
\end{table}

\subsection{Daltonizm}
\MyFigure[0.5\textwidth]{images/schematy/dalton.png}{Schemat koloru dla daltonistów}{fig:dalton}{\parencite{makieta}}
Schemat kolorów oparty na palecie Bank Wong dla osób z daltonizmem.

\begin{table}[ht]
\centering
\caption{Schemat kolorów dla Daltonizm}
\label{tab:dalt}
\begin{tabular}{|>{\raggedright\arraybackslash}p{4cm}|c|c|}
\hline
\textbf{Element} & \textbf{Kolor Hex} & \textbf{Kontrast} \\
\hline
funkcja       & \#56B4E9 & 7.22:1 \\
operator      & \#CCCCCC & 10.38:1 \\
słowo kluczowe & \#E69F00 & 7.4:1 \\
zmienna       & \#80DDE0 & 10.58:1 \\
klasa         & \#00C792 & 7.6:1 \\
typ           & \#B8A2EC & 7.47:1 \\
string        & \#F0E442 & 12.6:1 \\
stała         & \#D897BB & 7.19:1 \\
liczba        & \#F4935A & 7.27:1 \\
numer linii   & \#AED8EF & 11:1 \\
\hline
ostrzeżenie   & \#FFD04D & 11.41:1 \\
error         & \#63D5D9 & 9.56:1 \\
\hline
\end{tabular}
\newline\newline
{\raggedright\fontsize{10pt}{12pt}\selectfont\textit{Źródło: Opracowanie własne}\par}
\end{table}

\subsection{Pusty}
\MyFigure[0.5\textwidth]{images/schematy/blank.png}{Schemat koloru pusty}{fig:blank}{\parencite{makieta}}
Cały tekst edytora jest biały dla maksymalnego kontrastu. Dla osób z ostrymi problemami z widzeniem.

\begin{table}[ht]
\centering
\caption{Schemat kolorów pusty}
\label{tab:bialy}
\begin{tabular}{|>{\raggedright\arraybackslash}p{4cm}|c|c|}
\hline
\textbf{Element} & \textbf{Kolor Hex} & \textbf{Kontrast} \\
\hline
wszystkie elementy & \#FFFFFF & 16.67:1 \\
\hline
\end{tabular}
\newline\newline
{\raggedright\fontsize{10pt}{12pt}\selectfont\textit{Źródło: Opracowanie własne}\par}
\end{table}